% =============================================================================
%  §0a  先建直觉：OT in Plain Words
%  非数学背景读者的零门槛 on-ramp；置于正式预备节之前。
% =============================================================================
\TLsection{先建直觉 · OT in Plain Words（写给非数学背景读者）}

% ── 一句话直觉 ───────────────────────────────────────────────────────────────
\begin{frame}{一句话：最优传输在搬什么}
  你有一堆“东西”，形状像 $\mu$；你想把它重新摆成形状 $\nu$。
  从位置 $x$ 搬一单位到位置 $y$ 要付成本 $c(x,y)$，通常就是距离或距离的平方。
  \textbf{OT} 问的是：所有搬法里，哪一种总成本最低？

  \begin{center}
  \begin{tikzpicture}[x=0.74cm,y=0.50cm,>={Stealth[length=5pt]},
      barL/.style={fill=TLprimary,draw=TLprimaryDark},
      barR/.style={fill=TLaccent,draw=TLaccent},
      lab/.style={font=\scriptsize,text=TLinkSoft},
      arr/.style={->,TLink,thick}]
    \draw[TLinkSoft] (-0.2,0)--(2.6,0);
    \draw[TLinkSoft] (7.2,0)--(10.0,0);
    \foreach \x/\h in {0/2.2,0.8/3.5,1.6/1.5}
      \draw[barL] (\x,0) rectangle ++(0.55,\h);
    \foreach \x/\h in {7.5/1.4,8.3/2.5,9.1/3.3}
      \draw[barR] (\x,0) rectangle ++(0.55,\h);

    \node[lab] at (1.05,-0.55) {源 $\mu$};
    \node[lab] at (8.55,-0.55) {目标 $\nu$};
    \node[font=\small\bfseries,text=TLaccent] at (4.95,4.15) {最优传输};

    \draw[arr] (0.28,2.30) .. controls (3.0,3.2) and (5.6,3.1) .. (7.78,1.48);
    \draw[arr] (1.08,3.58) .. controls (3.1,4.0) and (5.9,3.7) .. (8.58,2.58);
    \draw[arr] (1.88,1.60) .. controls (3.3,1.2) and (6.2,1.7) .. (9.38,3.38);
  \end{tikzpicture}\par
  {\scriptsize 示意图。}
  \end{center}

  \TLtakeaway{最优传输把“从哪来、到哪去、花多少钱”放在同一个问题里：在所有可行搬法中找总成本最小者。}
\end{frame}

% ── 口算例子 ───────────────────────────────────────────────────────────────
\begin{frame}{一个能口算的例子}
  设 $\mu$ 放在 $\{0,1,2\}$ 上且每点权重为 $\tfrac13$，
  $\nu$ 放在 $\{1,3,5\}$ 上且每点权重为 $\tfrac13$；成本为 $c(x,y)=|x-y|$。
  顺序对齐给出最优方案，交叉搬法会更贵：
  \[
  \begin{aligned}
    0\to1,\quad 1\to3,\quad 2\to5:
      &\quad \frac{|0-1|+|1-3|+|2-5|}{3}
       =\frac{1+2+3}{3}=2,\\
    0\to5,\quad 1\to3,\quad 2\to1:
      &\quad \frac{|0-5|+|1-3|+|2-1|}{3}
       =\frac{5+2+1}{3}=\frac{8}{3}>2.
  \end{aligned}
  \]
  第一行没有交叉，第二行把左端点送到最远处，白白多走路。

  \TLtakeaway{一维最优就是"从小到大对齐"（单调匹配）。}
\end{frame}

% ── Monge vs Kantorovich ───────────────────────────────────────────────────
\begin{frame}{两种搬法：整粒搬 vs 可分流}
  \begin{columns}[T]
    \begin{column}{0.50\textwidth}
      \begin{itemize}\small
        \item \textbf{Monge：} 每个源点把自己的量整粒送到一个目标；这是一张映射 $T$。它可能没有解，因为一个点不能同时送到两个地方。
        \item \textbf{Kantorovich：} 允许把量分到多个目标；用流量表 $(\pi_{ij})$ 表示，行和是供应，列和是需求。这样总有解。
      \end{itemize}
    \end{column}
    \begin{column}{0.50\textwidth}
      \centering
      \begin{tikzpicture}[x=0.9cm,y=0.55cm,>={Stealth[length=5pt]},
          cell/.style={draw=TLinkSoft,minimum width=0.95cm,minimum height=0.45cm,
                       align=center,font=\scriptsize},
          head/.style={font=\scriptsize\bfseries,text=TLprimary},
          sumlab/.style={font=\scriptsize,text=TLinkSoft}]
        \node[head] at (1.55,3.35) {目标};
        \foreach \j/\lab in {1/$y_1$,2/$y_2$,3/$y_3$}
          \node[head] at (\j,2.8) {\lab};
        \node[head,rotate=90] at (-0.25,1.65) {源};
        \foreach \i/\lab in {1/$x_1$,2/$x_2$}
          \node[head] at (0,2.5-\i*0.55) {\lab};

        \node[cell,fill=TLprimaryTint] at (1,1.95) {$\tfrac12$};
        \node[cell,fill=TLprimaryTint] at (2,1.95) {$\tfrac12$};
        \node[cell] at (3,1.95) {$0$};
        \node[cell] at (1,1.40) {$0$};
        \node[cell,fill=TLprimaryTint] at (2,1.40) {$\tfrac12$};
        \node[cell,fill=TLprimaryTint] at (3,1.40) {$\tfrac12$};

        \node[sumlab] at (4.15,1.95) {行和 $1$};
        \node[sumlab] at (4.15,1.40) {行和 $1$};
        \node[sumlab] at (1,0.75) {列和 $\tfrac12$};
        \node[sumlab] at (2,0.75) {列和 $1$};
        \node[sumlab] at (3,0.75) {列和 $\tfrac12$};
        \draw[->,TLaccent,thick] (0.72,2.25) -- (1.18,2.10);
        \draw[->,TLaccent,thick] (0.72,2.25) -- (2.00,2.10);
        \node[font=\scriptsize,text=TLaccent] at (2.05,0.22) {一个源可分流};
      \end{tikzpicture}\par
      {\scriptsize 示意图。}
    \end{column}
  \end{columns}

  \TLtakeaway{本讲义用 Kantorovich（流量表）——总有解、且对 $\pi$ 线性。}
\end{frame}

% ── 对偶价格直觉 ───────────────────────────────────────────────────────────
\begin{frame}{另一个视角：价格（对偶）}
  想象一个中介在源位置 $x$ 以价格 $\psi(x)$ 收货，再在目标位置 $y$ 以价格 $\phi(y)$ 卖出。
  诚实定价的意思是：一次转手赚到的钱，不能超过自己真的把货从 $x$ 搬到 $y$ 的成本：
  \[
    \phi(y)-\psi(x)\le c(x,y).
  \]
  在这个限制下，中介能赚到的最大总利润，恰好等于你亲自安排搬运时必须支付的最小总成本。

  \TLtakeaway{"搬多少钱"="值多少钱" 在最优处相等——这就是 §2 的 Kantorovich 对偶。}
\end{frame}

% ── 应用地图 ───────────────────────────────────────────────────────────────
\begin{frame}{OT 在你的领域里}
  \begin{columns}[T]
    \begin{column}{0.64\textwidth}
      \begin{itemize}\small
        \item \textbf{机器学习：} Wasserstein 距离 / WGAN——衡量两个分布有多"远"；域自适应。
        \item \textbf{统计：} 比较分布、拟合优度检验。
        \item \textbf{图像处理：} 颜色迁移、形状插值（变形）。
        \item \textbf{经济学：} 匹配市场（Monge--Kantorovich 均衡）。
        \item \textbf{生物信息：} 单细胞发育轨迹重建。
      \end{itemize}
      这份讲义把这些应用共同依赖的数学地基讲全。
    \end{column}
    \begin{column}{0.36\textwidth}
      \centering
      \begin{tikzpicture}[x=0.55cm,y=0.55cm,>={Stealth[length=5pt]},
          node/.style={draw=TLprimary,fill=TLprimaryTint,rounded corners=3pt,
                       inner sep=3pt,font=\scriptsize,align=center},
          spoke/.style={->,TLink,thick}]
        \node[draw=TLaccent,fill=TLaccentLight,rounded corners=14pt,
              inner sep=5pt,font=\small\bfseries,text=TLaccent] (ot) at (0,0) {OT};
        \node[node] (ml) at (-2.3,1.7) {机器\\学习};
        \node[node] (st) at (2.2,1.7) {统计};
        \node[node] (im) at (-2.5,-0.2) {图像\\处理};
        \node[node] (ec) at (2.4,-0.2) {经济学};
        \node[node] (bio) at (0,-2.1) {生物\\信息};
        \foreach \n in {ml,st,im,ec,bio}
          \draw[spoke] (ot) -- (\n);
      \end{tikzpicture}\par
      {\scriptsize 示意图。}
    \end{column}
  \end{columns}

  \TLtakeaway{应用名字不同，核心问题相同：怎样比较、匹配、移动两个形状。}
\end{frame}

% ── 阅读路径 ───────────────────────────────────────────────────────────────
\begin{frame}{怎么读这份讲义（阅读路径）}
  \begin{itemize}\small
    \item ★ \textbf{必读}：动机/直觉帧、加粗标出的\emph{定义 / 定理}帧、可算例子、每帧的\emph{要点}。
    \item ⏭ \textbf{第一遍可跳}：标题含"证明 / 证明思路"的帧（想深究再回看；§6 末另有更硬的"技术备用"）。
    \item \textbf{背景}：只需微积分 + 基本概率 + 线性代数；生僻术语在 §0 现学现用，每个首次出现的词都配"一句话直觉"。
    \item \textbf{记号不熟}：随时翻文末"术语表 / 记号表"。
  \end{itemize}

  \begin{center}
  \begin{tikzpicture}[x=1cm,y=0.8cm,>={Stealth[length=5pt]},
      readnode/.style={draw=TLprimary,fill=TLprimaryTint,rounded corners=3pt,
                       inner sep=4pt,font=\scriptsize\bfseries,text=TLink},
      arr/.style={->,TLaccent,thick}]
    \node[readnode] (a) at (0,0) {直觉};
    \node[readnode] (b) at (2.4,0) {定义};
    \node[readnode] (c) at (4.8,0) {例子};
    \node[readnode] (d) at (7.2,0) {要点};
    \draw[arr] (a)--(b);
    \draw[arr] (b)--(c);
    \draw[arr] (c)--(d);
  \end{tikzpicture}
  \end{center}

  \TLtakeaway{第一遍走"直觉 → 定义 → 例子 → 要点"，证明留到第二遍。}
\end{frame}
